% Chapter 2

\chapter{Literature Review}

\label{Chapter2}

\lhead{Chapter 2. \emph{Literature Review}}

This chapter looks at the four main areas that this thesis is based on: how well the PCIe protocol works, designing DMA architectures using AXI, ways to reduce overhead from systems and networking research, and modeling latency in on-chip interconnects using RTL and queueing theory. Then, in section 2, we identify the specific problem that this thesis aims to solve.

%----------------------------------------------------------------------------------------

\section{PCIe Protocol and Performance Characterisation}
\label{sec:ch2_pcie}

The \ac{PCIe} Base Specification \cite{pcisig2017base} defines a layered architecture -- transaction, data link, and physical layers -- in which application data is carried inside \acp{TLP}. Every \ac{TLP} carries a fixed-size header (typically 12--16 bytes for memory requests), is subject to credit-based flow control between the transmitter and receiver, and its delivery is acknowledged at the data-link layer. Budruk~et~al.~\cite{budruk2004pcie} provide a detailed treatment of this layering and of the round-trip credit and acknowledgement mechanisms that contribute fixed, size-independent latency to every transaction, independent of the physical-layer transfer rate.

Here's a rewritten version of the input text in a more human-like tone, mimicking the style and structure of the provided human samples: When it comes to understanding how fixed-overhead behaviour works in real systems, there's surprisingly little research out there. One study by Neugebauer et al. did look at how PCIe performs in end-host networking on regular servers. They found that when it comes to small DMA transfers - the kind that's most relevant to packet-processing NICs - the throughput is actually determined by the number of separate PCIe transactions, rather than the total amount of data being transferred. To get close to the link bandwidth, they had to batch multiple descriptors or payloads per transaction. Now, their measurements were taken on actual production hardware, and they treated the NIC and host as a single unit - a black box, essentially. But what they didn't do was create a detailed model of the datapath that would show exactly how much of the overhead is due to the PCIe transaction layer, and how much is due to the DMA engine's descriptor processing. They also didn't provide a way to experimentally adjust the balance between these two factors. This is what motivated us to break down the system into smaller modules, as we'll see in Chapter 4 of this thesis.

%----------------------------------------------------------------------------------------

\section{AXI-Based DMA Architectures}
\label{sec:ch2_axi_dma}

On the fabric side of a PCIe endpoint, the \ac{AMBA} \ac{AXI4} and \ac{AXI4-Stream} protocols \cite{arm2021axi, arm2010axistream} are the de facto standard for connecting a DMA engine to on-chip processing logic. AXI4 defines a memory-mapped, burst-oriented protocol with independent address and data channels, while AXI4-Stream strips away addressing entirely in favour of a simple, back-pressured (\texttt{tvalid}/\texttt{tready}) beat-by-beat handshake with sideband fields (\texttt{tlast}, \texttt{tkeep}, \texttt{tuser}, \texttt{tid}) that this thesis reuses directly to carry per-packet size and identity information through the pipeline (\autoref{sec:ch4_sideband}).

Commercial DMA engines, like the AMD-Xilinx AXI DMA and the DMA/bridge logic in the UltraScale+ PCIe endpoint block, use a descriptor-ring model to operate. Here's how it works: the host or an on-chip processor writes a descriptor with details like source address, destination address, and length to a ring buffer in memory. Then, it signals the engine using a doorbell register or interrupt. The engine has to fetch and parse this descriptor, which is a memory or PCIe transaction with its own latency, before it can start moving the payload. The product guides for these cores provide information on register maps, supported burst lengths, and peak throughput figures. However, they don't model descriptor-fetch latency as a function of transfer size, treating it as an implementation detail instead. In contrast, this thesis's DMA model makes the fixed descriptor-processing cost a clear, parameterized quantity, called SETUP CYCLES. This allows us to isolate and measure its contribution to small-packet latency. By doing so, we can better understand the impact of this cost on the overall performance of the DMA engine. For instance, when transferring small packets, the descriptor-fetch latency can become a significant bottleneck. By explicitly modeling this cost, we can identify areas for optimization and improve the overall efficiency of the DMA engine. This is particularly important in applications where low latency is crucial, such as in real-time systems or high-performance computing. In summary, the descriptor-ring model used by commercial DMA engines has its own set of challenges, particularly when it comes to descriptor-fetch latency. By making this cost explicit and parameterized, we can gain a better understanding of its impact on performance and identify opportunities for improvement. This can lead to more efficient and effective DMA engines, which are essential in a wide range of applications.

%----------------------------------------------------------------------------------------

\section{Overhead Amortisation and Batching Techniques}
\label{sec:ch2_batching}

The idea behind the batching unit in this thesis is simple: combining many small events into fewer, larger ones can make them more efficient. This concept has been around for a while in systems and networking research, but it's usually applied in software or at a higher level, not directly in hardware like an RTL DMA datapath. For example, Mogul and Ramakrishnan found that when a network receives a lot of packets, handling each one with an individual interrupt can cause problems and slow things down. They showed that by polling and combining interrupts, so that multiple packets are handled at once, the system can recover and work more efficiently. This is because the fixed cost of handling each interrupt is spread out over many packets, making it more cost-effective. Later, network interface controllers started including hardware timers that could delay interrupts until a certain number of packets had arrived or a timeout had expired. This added a bit of latency to each packet, but it greatly reduced the overhead of handling each packet individually. By trading off a small amount of extra delay for a big reduction in overhead, these controllers could handle more packets and improve overall performance. This approach has been used in various forms, but the key idea is the same: by combining small events into larger ones, we can make systems more efficient and improve their performance. In the context of this thesis, the batching unit applies this principle to the RTL DMA datapath, aiming to reduce overhead and improve efficiency in a similar way.

These methods work at the point where the hardware and the operating system meet, combining signals that interrupt the system rather than the actual movement of data. The concept of using the same principle a step further down - combining many small packets of data into one big frame before they get to the parts of the system that add a fixed cost, and then measuring how much this improves things directly in a simulation - is what this project is about. It puts this idea into action with a part called the batching unit and checks how well it works with real numbers.

%----------------------------------------------------------------------------------------

\section{Latency Modeling of Interconnects}
\label{sec:ch2_latency_modeling}

Analytically, the behaviour this thesis is concerned with -- a service time made up of a fixed component plus a component proportional to the size of the job being served -- is a standard construction in queueing theory. Kleinrock's treatment of queueing systems \cite{kleinrock1975queueing} formalises service-time models of exactly this fixed-plus-proportional form and their consequences for waiting time and utilisation, and provides the conceptual basis for the closed-form PCIe and DMA delay expressions derived in \autoref{sec:ch3_latency_model}. However, classical queueing analysis models the service discipline abstractly; it does not, by itself, validate whether a specific RTL implementation of a specific protocol actually exhibits the assumed fixed-plus-proportional behaviour, or by how much a specific mitigation (such as batching) improves it for a specific packet-size distribution. Closing that gap between an analytical model and a cycle-accurate, simulatable hardware implementation -- built using standard SystemVerilog RTL practice \cite{ieee2017systemverilog} and simulated with an industrial simulator \cite{xilinx2022vivadosim} -- is the methodological contribution of this thesis.

%----------------------------------------------------------------------------------------

\section{Research Gap}
\label{sec:ch2_gap}

Synthesising the four areas above, three gaps recur:

\begin{enumerate}
\item \textbf{No open, parameterised, per-stage latency model.} PCIe performance studies \cite{neugebauer2018pcie} characterise small-transfer inefficiency empirically on commodity hardware, but as a single opaque number rather than a model that separately attributes latency to the PCIe transaction layer versus the DMA descriptor-processing stage. Vendor IP documentation \cite{xilinx2021axidma, xilinx2021pcieintblock} specifies the mechanisms that create this overhead (TLP framing, descriptor fetch) but does not quantify it as a function of packet size.
* **No use of batching at the RTL level for PCIe-AXI DMA transfers**. The idea of coalescing is already used for network interrupts, as seen in previous work, but applying this concept directly to an AXI-Stream DMA pipeline is not well documented. This means merging small packets into a larger frame before the PCIe/DMA overhead, and then measuring how this affects latency at the cycle level, is not something you can find in the literature or vendor documentation.
* There's no proven connection between the theoretical model and actual hardware measurements. The fixed-plus-proportional service models are commonly used in queueing theory, as seen in Kleinrock's work, but we need to know if a real PCIe-AXI DMA RTL pipeline really works this way. We also want to understand how much a solution like batching can improve it. To do this, we need a version of the pipeline that can be measured and simulated, which currently doesn't exist in the available research. This is a problem because it means we can't be sure if our theories are correct or how well our solutions will work in practice.
\end{enumerate}

This research tackles all three gaps at once. It starts by creating a detailed model of a PCIe-AXI DMA pipeline in Chapters 3 and 4, which takes into account the size of packets and how that affects latency. Then, in Chapter 5, it puts the coalescing principle into practice by adding a batching unit directly into the design. Finally, in Chapters 6 and 7, it checks how well the model works and measures the improvement that batching brings by using precise measurements from a hardware latency monitor built into the simulated design.